기룡이는 새 학기를 맞이하여 일정표(To Do List)를 작성해 알찬 학교생활을 보내려고 한다.

기룡이는 매주 일주일 단위의 큰 계획을 세우고, 이를 다시 하루 단위로 나누고, 마지막으로 시간대별 계획까지 세분화하여 각 계획을 아래와 같은 형식으로 작성하였다.
$$\texttt{DAY HH:MM HH:MM SUBJECT}$$

\begin{itemize}
\item \texttt{DAY}는 요일을 의미한다.
\item \texttt{HH:MM}은 24시간 형식의 시각을 의미한다. 첫 번째 시각은 계획의 시작 시각, 두 번째 시각은 종료 시각이다. 
\item \texttt{SUBJECT}는 해당 시간에 공부할 내용을 의미한다.
\end{itemize}

월요일부터 일요일까지의 계획을 모두 작성한 기룡이를 지켜보던 솔리는 기룡이의 갓생을 방해하기 위해 특정 시점마다 놀자고 꼬드기기 시작했다. 

솔리가 기룡이를 꼬드긴 여러 시점이 주어질 때, 각 시점에 기룡이가 어떤 공부를 하고 있는지 확인하고, 만약 계획이 없다면 신나게 놀아보도록 하자. 시점이 시작 시각 또는 종료 시각과 정확히 일치하는 경우도 계획이 진행 중인 것으로 간주한다.